\documentclass[12pt]{article}
\usepackage[utf8]{inputenc}
\usepackage[T1]{fontenc}
\usepackage{hyperref}
\title{Old French Language Resources for K-12 Curriculum Development}
\author{Compiled by Claude Sonnet 5.5 (claude-sonnet-5-5) for Andr\'as Kornai}
\date{}
\begin{document}
\maketitle

\section{Introduction}
Old French (842--ca.\ 1400) has no native speakers. Search turned up scholarly resources but no school curricula or classroom materials.

\section{Primary Educational Resources}
No primers, textbooks or school programmes were found in search.

\section{Linguistic and Reference Materials}
The \emph{Dictionnaire \'etymologique de l'ancien fran\c{c}ais} (DEAF), Heidelberg: \url{https://en.wikipedia.org/wiki/Dictionnaire_%C3%A9tymologique_de_l%27ancien_fran%C3%A7ais}.

A list of online early Romance resources (University of British Columbia): \url{https://blogs.ubc.ca/earlyromance/online-resources/}.

\section{Religious and Community Resources}
None found in search.

\section{Digital and Audio Resources}
The Base de Fran\c{c}ais M\'edi\'eval (BFM), a lemmatized and tagged corpus of Old and Middle French texts hosted at ENS de Lyon: \url{https://www.tei-c.org/activities/projects/base-de-francais-medieval-old-french-corpus/}.

\section{Curriculum Development Context}
Resources found are aimed at researchers. A heritage curriculum would need to adapt corpus texts and dictionaries; no existing school programme was found.

\end{document}